\section{Introduction}
Molecular property prediction is a fundamental pillar of modern drug discovery. Specifically, optimizing the Absorption, Distribution, Metabolism, Excretion, and Toxicity (ADMET) profile of a drug candidate is critical, as unfavorable ADMET characteristics are a primary cause of late-stage clinical trial failures. With the advent of deep representation learning, large-scale foundation models like MoLFormer \cite{ross2022large} have emerged, providing dense, continuous, and context-aware embeddings of molecular structures. 

However, evaluating these embeddings in practical drug discovery pipelines presents a significant challenge. Unlike the massive corpora used to pre-train foundation models, experimental ADMET data is time-consuming and expensive to acquire. To standardize the evaluation of machine learning models in this domain, the Therapeutics Data Commons (TDC) introduced the ADMET Benchmark Group \cite{huang2021therapeutics}. This benchmark formalizes the pervasive "Curse of Dimensionality" ($P \gg N$) inherent to cheminformatics, as target datasets frequently contain fewer than a thousand unique samples.

In this work, we focus specifically on the \texttt{caco2\_wang} regression task from the TDC ADMET group. Caco-2 is a human colon epithelial cancer cell line widely utilized in pharmaceutical research as an \textit{in vitro} proxy to predict human intestinal absorption. Accurate modeling of Caco-2 permeability is essential for the design of orally administered drugs. However, mapping 768-dimensional MoLFormer embeddings to the $\approx 600$ training samples of the Caco-2 dataset creates an extreme low-sample environment that pushes standard predictive models to their theoretical limits.

Furthermore, in clinical and pharmaceutical deployments, high point accuracy is insufficient. When models encounter Out-of-Distribution (OOD) samples—molecules containing novel scaffolds or functional groups that lie far outside the training manifold—they must possess robust epistemic uncertainty quantification. A model must "know what it does not know"; flagging highly uncertain predictions prevents confident but catastrophic failures and allows for targeted human intervention. 

To address these challenges, we benchmark Stochastic Variational Deep Kernel Learning (SV-DKL) against classical XGBoost Deep Ensembles to assess their viability in the $P \gg N$ regime. We utilize the Optuna framework to systematically tune both architectures. Our contributions are as follows: first, we demonstrate that tree-based ensembles suffer from a fundamental geometric mismatch when partitioning dense continuous representations, leading to severe uncertainty collapse. Second, we show that SV-DKL, when stabilized with Spectral Normalization to prevent feature collapse, effectively bypasses this dimensionality bottleneck. By natively exploiting the spatial geometry of the foundation model embeddings, SV-DKL provides superior predictive accuracy alongside mathematically rigorous OOD uncertainty calibration.
